\SourceLeaf{586}\pdfbookmark[1]{Astronomische Nachrichten Nr. 417. 1841. Februar 9.}{nachrichten-417}
\begin{center}\small\textsc{Schumacher}.\quad Astronomische Nachrichten Nr.\ 417.\quad 1841.\quad Februar 9.\end{center}

In einer Zeitung las ich kürzlich die Nachricht, dass der amerikanische Marinekapitän \textsc{Wilkes}
\begin{center}\textit{dem magnetischen Südpol}\end{center}
ziemlich nahe gekommen sei und dass er auf $67^0\;4'$ südlicher Breite und $147^0\;30'$ Länge (zweifellos östlich von Greenwich) eine magnetische Abweichung von $12^0\;35'$ nach Osten und eine Neigung von $87^0\;30'$ gefunden habe. Nach einer rasch angestellten Rechnung würde ich den tatsächlichen Pol demnach vorläufig
\begin{center}auf $70^0\;21'$ südlicher Breite,\qquad $146^0\;17'$ Länge\end{center}
ansetzen. Dieser Ort liegt demjenigen, den meine Theorie [S.\ 163 d.\ B.] angegeben hat, viel näher, als ich selbst erwartet hatte. Der tatsächliche Pol, wie ich dort vermutet hatte, weiter nördlich als der nach der Theorie berechnete. Der Unterschied in der Breite erreicht aber nur ein Drittel dessen, worauf ich nach Durchsicht der Beobachtungen von Hobarttown gefasst war. Ebenso liegt der tatsächliche Pol weiter westlich als der nach der Theorie berechnete, und hier ist der Unterschied fast genau so groß wie der, den ich an der genannten Stelle angenommen hatte. Im Übrigen ist es unnötig zu erwähnen, dass in diesen hohen Breiten ein Unterschied von sechs Längengraden nur ebenso viel bedeutet wie zwei Breitengrade.

